\section{The mechanism works}
\label{sec:mechanism}

\begin{table}[t]
\centering
\caption{Signed deviation from the real radial spectrum in decibels; negative
means too little energy. The static ablation is indistinguishable from MSE.}
\label{tab:bands}
\begin{tabular}{lrrr}
\toprule
Configuration & Low & Mid & High \\
\midrule
MSE & $-0.36$ & $-0.80$ & $-2.20$ \\
Static wavelet ($p{=}0$) & $-0.35$ & $-0.79$ & $-2.23$ \\
Frequency-aware & $-0.31$ & $-0.56$ & $\mathbf{-1.79}$ \\
\bottomrule
\end{tabular}
\end{table}

\begin{figure}[t]
\centering
\includegraphics[width=0.85\linewidth]{spectrum.png}
\caption{Deviation from the real spectrum against normalised frequency. The
MSE and static-wavelet curves coincide across the entire axis: the wavelet
decomposition on its own changes nothing. Only the time-dependent schedule
separates from them, and the gap widens with frequency.}
\label{fig:spectrum}
\end{figure}

Table~\ref{tab:bands} and Figure~\ref{fig:spectrum} establish three things.

The premise holds: the deficit is roughly sixfold worse at high frequency than
at low. The objective's improvement is largest there in absolute terms
(\SI{0.41}{\decibel} against \SI{0.05}{\decibel} at low frequency), so the
correction lands where the method says it should. We state this precisely: the
improvement is \emph{broadband with a high-frequency tilt}, not an exclusively
high-frequency correction --- the proportional gain is largest in the mid band,
and a referee will recompute.

The wavelet decomposition contributes nothing. The static ablation, which keeps
the decomposition and removes only the dependence on $\sigma$, overlaps MSE
across the whole frequency axis. The temporal schedule is the entire effect.
The correction also survives the protocol of \S\ref{sec:confound}
($p = 0.0056$), so it is a genuine frequency effect rather than the timestep
confound in disguise.

\begin{figure}[t]
\centering
\includegraphics[width=\linewidth]{compare.png}
\caption{Samples from identical initial noise, one row per objective. A
\SI{0.4}{\decibel} high-frequency change at $32\times32$ is below the visible
threshold.}
\label{fig:compare}
\end{figure}

Finally, the correction is not visible. Figure~\ref{fig:compare} pairs samples
on their initial latent, so any difference would be attributable to the
objective; none is apparent.
